\documentclass{article}
\usepackage{graphicx} % Required for inserting images

\title{FIM 548 HW4 Writeup}
\author{Jacob Attia, Clay Ewing, Zach Wiebe}
\date{May 2026}

\usepackage{amsmath, amssymb}
\usepackage{geometry}
\usepackage{booktabs}
\usepackage{parskip}
\geometry{margin=1in}
\begin{document}


\maketitle
\subsection*{Problem 1: American Put Option Pricing}
\subsection*{Setup}

We price an American put option on a risky asset following GBM under the risk-neutral measure:
\[
dS_t = S_t(r\,dt + \sigma\,dB_t), \quad S_0 = 100
\]
with parameters $r = 0.05$, $\sigma = 0.4$, strike $K = 100$, maturity $T = 1$, discretized into $N = 5000$ time steps and $M = 200{,}000$ sample paths.

From Lecture 12, the fair price of an American option is given by the optimal stopping problem:
\[
V(0, S_0) = \sup_{\tau \in \mathcal{T}} \mathbb{E}\left[e^{-r\tau}(K - S_\tau)^+\right]
\]
where $\mathcal{T}$ is the set of all stopping times adapted to the filtration generated by $S$.

% ─────────────────────────────────────────────────────────────────────────────
\subsection*{Part 1: Martingale Duality Upper Bound}

From Lecture 13, the martingale duality result states that for any martingale $(M_t)$ with $M_0 = 0$:
\[
V(0, S_0) \leq \mathbb{E}\left[\sup_{t \in \{t_0, \ldots, t_N\}} \left(G_t - M_t\right)\right]
\]
where $G_t = e^{-rt}(K - S_t)^+$ is the discounted payoff process. Since the left-hand side does not depend on $M$, we have:
\[
V(0, S_0) \leq \inf_{\substack{M:\, M_0=0 \\ M \text{ martingale}}} \mathbb{E}\left[\sup_t (G_t - M_t)\right]
\]
and in fact equality holds for the optimal martingale $M^*$, which is the martingale part of the Doob decomposition of the Snell envelope $V(t, S_t)$. The closer $M$ is to $M^*$, the tighter the upper bound.

We test two candidate martingales:

\textbf{Martingale (a):}
\[
M_t^{(a)} = e^{-rt}S_t - S_0
\]


\textbf{Martingale (b):} 
\[
M_t^{(b)} = e^{-rt}P(t, S_t) - P(0, S_0)
\]
where $P(t, x)$ is the Black-Scholes put price:
\[
P(t, x) = Ke^{-r(T-t)}\Phi(-d_2) - x\Phi(-d_1)
\]


\textbf{Execution.} To avoid storing the full path matrix (which ended up being unable to compute due to approx 7.5 GB of memory), we stream through time steps sequentially, maintaining only the current stock price vector and the running supremum for each martingale. At each step $i$ we compute $G_{t_i}$, $M_{t_i}^{(a)}$, $M_{t_i}^{(b)}$, update the running supremum, then step the stock price forward. This reduces memory usage to a viable solution.

\textbf{Results.}

\begin{center}
\begin{tabular}{lcc}
Martingale & Upper Bound & 95\% CI \\
\midrule
(a) Discounted stock & 53.1859 & (53.0432, 53.3287) \\
(b) European put     & 13.8551 & (13.8513, 13.8589) \\
\bottomrule
\end{tabular}
\end{center}

Martingale (b) produces a dramatically tighter upper bound. This is because the European put value process approximates the true Snell envelope well in the continuation region since exercising early is not optimal. Martingale (a) has no structural relationship to the put payoff: when the stock drops (making the put valuable), $e^{-rt}S_t - S_0$ becomes very negative, causing $G_t - M_t^{(a)}$ to be artificially large and inflating the supremum far above the true price.

% ─────────────────────────────────────────────────────────────────────────────
\subsection*{Part 2: Candidate Stopping Time and Lower Bound}

Any specific stopping time $\tau$ gives a lower bound since:
\[
\mathbb{E}[G_\tau] \leq V(0, S_0) = \sup_\tau \mathbb{E}[G_\tau]
\]
From Lecture 13, the candidate stopping time is defined as:
\[
\tau = \inf\{t \in [0,1] : V_0 + M_t \leq G_t\} \wedge T
\]
where $V_0$ is the upper bound from Part 1. Together, Parts 1 and 2 produce the interval:
\[
\left[\mathbb{E}[G_\tau],\, V_0\right]
\]
which brackets the true price $V(0, S_0)$.

\textbf{Execution.} Since the stopping time depends on $V_0$ from Part 1, a second pass through the paths is required, and we use the same seed to reproduce identical paths. We evaluate the exercise condition $G_t - M_t \geq V_0$ at each step, recording $G_\tau$ at the first triggered time. 

\textbf{Results.}

\begin{center}
\begin{tabular}{lccc}
Martingale & $\mathbb{E}[G_\tau]$ & $V_0$ & Interval width \\
\midrule
(a) Discounted stock & 13.2802 & 53.1859 & 39.91 \\
(b) European put     & 13.4489 & 13.8551 & 0.41  \\
\bottomrule
\end{tabular}
\end{center}

Martingale (b) produces an extremely tight interval of width 0.41, pinning the true American put price between 13.4489 and 13.8551. The interval from martingale (a) is practically useless due to the loose upper bound.

Both stopping times give $\mathbb{E}[G_\tau] > P_0 = 13.1459$, confirming that early exercise adds value relative to holding until maturity. 

% ─────────────────────────────────────────────────────────────────────────────
\subsection*{Part 3: Least-Squares Monte Carlo (LSMC)}

From Lecture 13, the LSMC (Longstaff-Schwartz) algorithm prices the American put by approximating the continuation value at each time step via regression. 
\[
V(i, s) = \max\left\{(K-s)^+,\, \mathbb{E}\left[e^{-r\Delta t}V(i+1, S_{i+1}) \mid S_i = s\right]\right\}
\]
We approximate the conditional expectation as a linear combination of basis functions:
\[
\mathbb{E}[V(i+1, S_{i+1}) \mid S_i = x] \approx \sum_{k=1}^{7} \beta_{ik}\psi_k(x), \qquad \psi_k(x) = x^{k-1}
\]
The coefficients are estimated by OLS regression of the discounted next-period values $e^{-r\Delta t}V(i+1, S_{i+1}^{(\ell)})$ on the basis matrix $(\psi_k(S_i^{(\ell)}))_{\ell,k}$, using QR decomposition for numerical stability.

Regression is performed only on in-the-money paths since it is never optimal to exercise out of the money.

\textbf{Execution.} Storing $M \times (N+1)$ paths simultaneously is not feasible with the instructed N = 5000, M = 200,000. LSMC is therefore run in 20 independent batches of 10,000 paths each. Each batch runs independently, producing a price estimate $\hat{V}^{(b)}$. The final price is the average:
\[
\hat{V} = \frac{1}{20}\sum_{b=1}^{20} \hat{V}^{(b)}
\]

\textbf{Results.}
\[
\text{LSMC price} = 13.6879, \qquad 95\%\ \text{CI}: (13.6249,\ 13.7509)
\]

The LSMC price falls within the duality interval from Part 2:
\[
\underbrace{13.4489}_{\text{lower bound}} \;\leq\; \underbrace{13.6879}_{\text{LSMC}} \;\leq\; \underbrace{13.8551}_{\text{upper bound}}
\]
confirming consistency across all three methods. The American put premium over the European put is $13.6879 - 13.1459 = 0.542$, reflecting the value of early exercise.

\end{document}

\end{document}
